\section{Background and Related Work}
\label{sec:background}

\paragraph{Standardized benchmarks.} MLPerf~\cite{reddi2020mlperf} anchors ML
performance measurement: its closed division fixes model architectures and
weights so that hardware can be compared fairly. The structural choices this
requires, however, create coverage gaps for systems research. MLPerf segments
by hardware scale --- \emph{Tiny}~\cite{banbury2021mlperftiny} for
microcontrollers, \emph{Edge}, \emph{Datacenter} --- leaving a mid-range gap:
devices such as Jetson-class boards, Apple-silicon workstations, and now
GB10-class unified-memory nodes are too capable for Tiny yet unrealistically
provisioned for Datacenter workloads. Researchers who want a \emph{continuous}
scalability axis --- where does this system break as the model grows? ---
cannot get it from discrete tracks. The broader ecosystem is fragmented:
EEMBC MLMark targets embedded inference but has not tracked the transformer
era; AI Benchmark~\cite{ignatov2019aibenchmark} covers mobile SoCs as a
consumer scorecard; TPCx-AI~\cite{brucke2023tpcxai} models end-to-end
enterprise pipelines but at minimum scale factors that exceed edge budgets,
over predominantly classical-ML workloads. AISBench~\cite{aisbench}, the
closest recent effort, argues for end-to-end multi-stream measurement on
large servers and clusters; \sysname{} shares its diagnosis but targets the
opposite end of the spectrum and adds declarative graph reconfiguration and
per-process collocation attribution as first-class capabilities.

\paragraph{The measurement boundary.} To isolate accelerator performance,
standard harnesses preload preprocessed samples into memory and time only the
model pass. This is the right control for hardware comparison and the wrong
one for systems analysis: data loading, decoding, and preprocessing routinely
rival model time in vision and multimodal
pipelines~\cite{coleman2017dawnbench}, and retrieval dominates many RAG
configurations. A benchmark that cannot see these stages cannot rank systems
by the metric deployments care about. \sysname{} makes every stage a
separately-timed, first-class citizen and can either preload (to reproduce
the MLPerf-style control) or perform just-in-time I/O (to measure the real
path) --- the choice is a configuration flag, not a fork of the harness.

\paragraph{Linear execution.} The prevailing benchmark abstraction is a
static feed-forward pass: a request enters, one network runs, a result
leaves. Compound AI systems break this: agentic loops introduce cyclic,
budget-bounded control flow; RAG pipelines branch on the semantics of
intermediate outputs~\cite{lewis2020rag,asai2024selfrag}; self-reflective
systems re-enter retrieval with rewritten queries. These are not exotic
research artifacts --- they are the dominant deployment pattern for LLM
applications --- and their systems behavior (per-stage residency, cross-query
overlap, retry-tail latency) is invisible to a fixed-graph harness.
\sysname{}'s pipelines are directed graphs with cycles, fan-out, and
configurable fan-in policies precisely so this class of workload can be
expressed, measured, and \emph{ablated}.

\paragraph{Collocation and attribution.} Benchmarks measure isolation;
deployments collocate. An edge node serving inference is simultaneously
fine-tuning, re-indexing its retrieval corpus, or hosting a second tenant.
On discrete-GPU servers, collocation contends chiefly for device time and
can be partitioned with MPS or MIG; on \emph{unified-memory} systems the
contended resource is the shared DRAM bandwidth itself, which every engine
--- CPU, GPU, neural accelerator --- draws from one pool. TPCx-AI's
concurrent streams measure aggregate throughput but attribute nothing;
MLPerf treats multi-tenancy as out of scope. \sysname{} runs each pipeline
as its own OS process with per-process resource listeners, so interference
is not merely induced but \emph{attributed}: which pipeline, which stage,
which engine, and --- with the memory-controller counters we exploit in
\S\ref{sec:results-staged} --- how many bytes.

\paragraph{Benchmarks as research instruments.} Finally, standardization and
exploration pull in opposite directions. Fair comparison demands monolithic,
tightly-controlled implementations; systems research demands the opposite ---
swap the retriever, halve the model, add a co-runner, and observe. In a
monolithic harness each such edit is an invasive, non-transferable code
change. In \sysname{} it is a configuration diff, and \S\ref{sec:results}
uses exactly this property as an experimental method: our contention study
descends from system-level observation to mechanism through a sequence of
single-element configuration changes, each shown explicitly.
